\documentclass[sigconf,screen]{acmart}

\usepackage{booktabs}
\usepackage{graphicx}
\usepackage{xcolor}
\usepackage{pifont}
\usepackage{enumitem}
\usepackage{subcaption}
\usepackage{tikz}
\usepackage{amsmath}
\usepackage{amssymb}
\usetikzlibrary{arrows.meta, positioning, shapes.geometric, calc, decorations.pathreplacing}

% Colors for framework visualization
\definecolor{awarenessblue}{HTML}{3B82F6}
\definecolor{negotiationamber}{HTML}{F59E0B}
\definecolor{bnagreen}{HTML}{10B981}
\definecolor{gategray}{HTML}{6B7280}
\definecolor{constraintpurple}{HTML}{8B5CF6}
\definecolor{perceptionlight}{HTML}{DBEAFE}
\definecolor{comprehensionlight}{HTML}{BFDBFE}
\definecolor{projectionlight}{HTML}{93C5FD}
\definecolor{proposelight}{HTML}{FEF3C7}
\definecolor{deliberatelight}{HTML}{FDE68A}
\definecolor{resolvelight}{HTML}{FCD34D}
\definecolor{loopbg}{HTML}{F8FAFC}

\newcommand{\system}{\textsc{GroundingKit}}
\newcommand{\todo}[1]{\textcolor{red}{[TODO: #1]}}
\newcommand{\dimicon}[1]{\tikz[baseline=-0.5ex]{\node[draw, rounded corners=2pt, fill=#1, minimum size=0.35cm] {};}}
\newcommand{\moment}[1]{\textsc{#1}}

\begin{document}

\title{Designing Coordination Infrastructure for Collaborative Writing}

\author{[Authors]}
\affiliation{\institution{[Institution]}}
\email{[email]}

\begin{abstract}
Collaborative writing requires teams to continuously negotiate shared understanding---what to write, how individual contributions relate to the collective goal, and when that goal should evolve. Yet current tools support only change-level awareness (track changes, real-time cursors, comments), leaving intent-level coordination unsupported. Teams can see \emph{what} each other wrote but cannot assess \emph{whether} that writing serves their shared plan or renegotiate the plan when it should change.

We address this gap in two steps. First, we define a \emph{coordination design space} for collaborative writing, organized around a boundary negotiating artifact (BNA) that externalizes team consensus. The BNA and writing form a \emph{negotiation space} connected by four types of relationships; a coordination pipeline---Awareness, Gate, Negotiation (Propose $\rightarrow$ Deliberate $\rightarrow$ Resolve)---operates on these relationships, with configurable dimensions at each stage. We derive this design space from a formative study with 3 teams (11 participants) and ground it in BNA theory, situational awareness, and conversation for action.

Second, we present \system{}, an open platform that operationalizes this design space. \system{} provides a foundation layer (abstract data model + view layer with entity-based UI slots), two declarative protocols (Awareness Protocol for sensing capabilities, Negotiation Protocol for coordination paths), and a guided team interface for configuring coordination pipelines. Developers contribute new capabilities through the same protocol used for built-in functions; teams configure their pipeline by answering questions about their collaboration context. We deploy \system{} with [N] teams and report findings on how teams define and use coordination infrastructure in practice.
\end{abstract}

\maketitle

\input{sections/introduction}
\input{sections/related-work}
\input{sections/formative-study}
\section{Design Space for Coordination in Collaborative Writing}
\label{sec:framework}

We propose to reconceive collaborative writing as a process of negotiation around team shared understanding. Current tools treat writing itself as the sole medium of collaboration---all coordination happens in the text stream, through track changes, comments, and ad hoc communication. We propose a different structure: teams maintain a shared understanding representation (BNA) independent of the writing itself; individual writing proceeds within the framework of this shared understanding; and when the shared understanding needs to change, the team negotiates through an explicit pipeline.

This section first defines BNA and the relationships between BNA and writing that form a \emph{negotiation space} (\S\ref{sec:bna}), then defines the coordination design space that operates on this negotiation space (\S\ref{sec:design-space}).


\subsection{BNA and the Negotiation Space}
\label{sec:bna}

\subsubsection{BNA in Collaborative Writing}

Lee's Boundary Negotiating Artifact (BNA)~\cite{lee2007boundary} describes a phenomenon in non-routine collaboration: when no pre-established standards exist, participants create artifacts that establish and evolve consensus---artifacts that do not passively carry information but actively participate in negotiation. Collaborative writing is exactly such a setting: every project requires teams to build and continuously maintain consensus about ``what we are writing.''

We apply BNA to collaborative writing, defining it as the externalized structure of team shared understanding---the consensus about ``what to write'' that the team collectively follows. With this structure, writing acts gain coordination meaning relative to team consensus; implicit cross-section relationships become trackable; negotiation can happen at the intent level rather than the sentence level; and awareness and negotiation have concrete anchors.

To serve these roles, a BNA must satisfy five conditions:

\begin{description}[leftmargin=1.5em, nosep, itemsep=3pt]
\item[C1: Composed of intent units.] The basic unit of BNA is an \emph{intent unit}---each expressing the team's consensus about what a part of the document ``should convey.'' Intent units are more abstract than specific writing, allowing different implementations while enabling negotiation at the intent level.

\item[C2: Has internal dependency structure.] Intent units have explicit dependencies (depends-on, supports, must-be-consistent). These dependencies are implicit in prose---BNA makes them visible, enabling cross-section impact tracking.

\item[C3: Has ownership.] Each intent unit has a designated owner, giving coordination a concrete addressee.

\item[C4: Traceable to writing.] BNA and writing maintain a traceable correspondence. At any moment, one can answer ``has this intent been realized?'' Without this traceability, BNA is a dead checklist.

\item[C5: Dynamically negotiable.] BNA is not a fixed plan but a continuously evolving shared understanding. Team members can propose modifications, but modifications require negotiation---this is why the pipeline exists.
\end{description}


\subsubsection{Concrete Forms of BNA}

BNA satisfying C1--C5 can take different concrete forms. Two examples illustrate how intent units, dependencies, and ownership look in different writing contexts.

\paragraph{Example 1: Living Outline.} A team co-writes an academic paper. Their BNA is a living outline where each entry is an intent unit:

\begin{table}[h]
\small
\begin{tabular}{@{}llp{5.5cm}l@{}}
\toprule
\textbf{ID} & \textbf{Intent} & \textbf{Description} & \textbf{Owner} \\
\midrule
I1 & Introduction & Motivate the coordination problem with a real scenario & Alice \\
I2 & Research Q. & State how to provide coordination infrastructure for CW & Alice \\
I3 & System Design & Justify the living outline as externalized consensus & Bob \\
I4 & Evaluation & Report deployment study methods and findings & Carol \\
\bottomrule
\end{tabular}
\end{table}

\noindent Dependencies: I3 \emph{depends-on} I2 (system design must respond to the research question); I4 \emph{must-be-consistent} I3 (evaluation must correspond to system design). If Alice revises the research question I2, Bob's system design I3 may need adjustment---a dependency invisible in prose but explicit in BNA.

\paragraph{Example 2: Argument Map.} A team co-writes a policy recommendation. Their BNA is an argument map where each node is an intent unit:

\begin{table}[h]
\small
\begin{tabular}{@{}llp{4.5cm}ll@{}}
\toprule
\textbf{ID} & \textbf{Intent} & \textbf{Description} & \textbf{Type} & \textbf{Owner} \\
\midrule
C1 & Main claim & AI systems need mandatory transparency & claim & Alice \\
E1 & Evidence & EU AI Act implementation data & evidence & Bob \\
C2 & Counter & Excessive transparency harms competitiveness & counter & Carol \\
E2 & Evidence & Tech company feedback surveys & evidence & Carol \\
\bottomrule
\end{tabular}
\end{table}

\noindent Dependencies: E1 \emph{supports} C1; C2 \emph{challenges} C1; E2 \emph{supports} C2. Both forms differ in structure but satisfy C1--C5.

% ─── FIGURE: Two BNA Forms ───
\begin{figure}[t]
\centering
\fbox{\parbox{0.95\columnwidth}{\centering\vspace{2cm}
\textbf{Figure: Two BNA Forms}\\[0.5em]
\small Living outline (left) and argument map (right), each showing intent units with dependency edges and ownership annotations. Both satisfy C1--C5 despite different structures.
\vspace{2cm}}}
\caption{Two concrete forms of BNA. A living outline (left) organizes intent units hierarchically with typed dependencies; an argument map (right) organizes them as claims, evidence, and counter-claims with support/challenge edges. Both externalize team consensus and make cross-section dependencies explicit.}
\label{fig:bna-forms}
\end{figure}


\subsubsection{The Negotiation Space}

BNA, Writing, and the relationships between them form the \textbf{negotiation space} of collaborative writing---the foundation on which coordination operates.

BNA is the team's shared understanding---consensus about ``what to write.'' Writing is the individual's realization of this consensus. Previously, both were entangled in the same document. With BNA, shared understanding has its own carrier, and writing returns to the individual's implementation space.

Four types of relationships exist in the negotiation space:

\begin{description}[leftmargin=1.5em, nosep, itemsep=3pt]
\item[BNA $\leftrightarrow$ BNA:] Dependencies between intent units. I3 \emph{depends-on} I2; E1 \emph{supports} C1. These make the document's structural logic explicit.

\item[BNA $\rightarrow$ Writing:] Shared understanding guides individual writing. Each intent unit defines expectations for the corresponding writing. Bob writes the system design following I3's intent; how he writes it is his freedom.

\item[Writing $\rightarrow$ BNA:] Writing feeds back to shared understanding. While writing, the relationship between content and intent continuously changes---drifting, discovering new angles, finding an intent no longer appropriate.

\item[Writing $\leftrightarrow$ Writing:] Indirect relationship through BNA. Alice's Introduction and Bob's Method appear independent in the text stream, but because I3 \emph{depends-on} I2, Alice's revision to the research question propagates to Bob's system design. This relationship is invisible in prose but trackable through BNA.
\end{description}

\noindent Changes in these relationships are coordination triggers. With trackable relationships, it becomes possible to build a coordination pipeline on top of them.

% ─── FIGURE: Negotiation Space ───
\begin{figure*}[t]
\centering
\fbox{\parbox{0.95\textwidth}{\centering\vspace{3cm}
\textbf{Figure: The Negotiation Space}\\[0.5em]
\small BNA (left, team shared understanding) and Writing (right, individual realization) connected through four types of relationships. (1) Internal dependencies between intent units. (2) Intent guides writing. (3) Writing feeds back to intent. (4) Writing sections are indirectly related through intent dependencies---invisible in text but explicit through BNA.
\vspace{3cm}}}
\caption{The negotiation space. BNA and Writing form two layers connected by four types of relationships. The BNA externalizes team consensus and makes the structural logic of the document explicit; Writing is the individual realization of this consensus. Changes in these relationships trigger coordination.}
\label{fig:negotiation-space}
\end{figure*}


\subsection{Coordination Design Space}
\label{sec:design-space}

Section~\ref{sec:bna} defined the negotiation space: BNA externalizes team consensus, Writing is the individual realization, and four types of relationships connect them. Current tools lack this space---track changes records who changed which words, but no one knows how those changes relate to team consensus.

BNA enables coordination through a core mechanism: it lets writers \emph{navigate between intent (abstract) and writing (concrete)}. A writer can start from their writing and understand its position relative to intent; or start from intent and project a modification into writing, seeing its concrete impact on themselves and others. This bidirectional navigation between abstract and concrete does not exist in scroll-bound text---it is the foundational capability BNA provides for coordination.

Based on this capability, we define the coordination design space (Figure~\ref{fig:pipeline}). Writers compose in personal space; \textbf{Awareness} continuously senses changes in the relationships between writing and BNA; when changes reach a team-configured threshold, the writer crosses the \textbf{Gate} into team space; \textbf{Negotiation} updates the BNA through Propose $\rightarrow$ Deliberate $\rightarrow$ Resolve; the updated BNA feeds back into writing and awareness---forming a continuous coordination loop.

% ─── FIGURE: Coordination Pipeline ───
\begin{figure*}[t]
\centering
\fbox{\parbox{0.95\textwidth}{\centering\vspace{3cm}
\textbf{Figure: Coordination Pipeline}\\[0.5em]
\small Top layer: BNA (team consensus) + Writing (individual). Loop: Writing (personal space) $\rightarrow$ Awareness (sense relationship changes) $\rightarrow$ Gate (threshold) $\rightarrow$ Negotiation (team space: Propose $\rightarrow$ Deliberate $\rightarrow$ Resolve) $\rightarrow$ Update BNA $\rightarrow$ back to Writing. Gate is a threshold condition between Awareness and Negotiation, not an independent stage. Awareness dimensions marked $\bigstar$ form the Gate condition.
\vspace{3cm}}}
\caption{The coordination pipeline. Awareness operates in personal space, continuously sensing the writer's relationship to team consensus. When $\bigstar$-marked dimensions (Scope, Severity, Group Interest) exceed a team-configured threshold, the Gate triggers entry into team space for Negotiation. The loop cycles continuously: writing changes relationships, awareness detects changes, negotiation updates BNA, updated BNA reshapes awareness.}
\label{fig:pipeline}
\end{figure*}


\subsubsection{Awareness}
\label{sec:awareness}

BNA externalizes the otherwise implicit relationships between writing and team consensus, enabling writers to perceive their position in the negotiation space. The Awareness design space has three groups of dimensions.

\paragraph{When --- triggering awareness.}

\begin{table}[h]
\small
\begin{tabular}{@{}lp{2.5cm}p{4cm}@{}}
\toprule
\textbf{Dimension} & \textbf{Design Question} & \textbf{Option Space} \\
\midrule
Trigger & What triggers awareness & Ambient, on-demand, post-action, external (others' changes) \\
Timing & Before or after writing & Pre-action (check alignment before writing), post-action (review drift after writing) \\
\bottomrule
\end{tabular}
\end{table}

\paragraph{What --- sensing relationship changes.}

\begin{table}[h]
\small
\begin{tabular}{@{}lp{2cm}p{4.5cm}@{}}
\toprule
\textbf{Dimension} & \textbf{Design Question} & \textbf{Option Space} \\
\midrule
Writing $\leftrightarrow$ Intent & My writing vs.\ my intent & Covered, partial, deviated, extended, orphan \\
Writing $\leftrightarrow$ BNA & My writing vs.\ team consensus & Aligned, in tension, diverged \\
Writing $\leftrightarrow$ Others & My writing vs.\ others' (via deps) & Unrelated, linked, potential conflict, conflict \\
Decision History & Past resolutions constrain? & Unconstrained, constrained by prior resolution \\
\bottomrule
\end{tabular}
\end{table}

\paragraph{Impact $\bigstar$ Gate --- assessing consequences.}

\begin{table}[h]
\small
\begin{tabular}{@{}lp{2cm}p{4.5cm}@{}}
\toprule
\textbf{Dimension} & \textbf{Design Question} & \textbf{Option Space} \\
\midrule
Scope $\bigstar$ & Who is affected & None, specific (via deps), team-wide \\
Severity $\bigstar$ & How significant & Cosmetic, local, directional, fundamental \\
Propagation & How it spreads & Only my intent, direct dependents, multi-level chain \\
Group Interest $\bigstar$ & Does team need to know & Need-to-know, want-to-know, won't care, unknown \\
\bottomrule
\end{tabular}
\end{table}

\noindent Dimensions marked $\bigstar$ jointly form the \textbf{Gate}: when their combination reaches a team-configured threshold, the writer transitions from personal space to team space. Different teams set different thresholds---a high-trust pair may trigger only on \emph{fundamental} + \emph{team-wide}; a high-stakes distributed team may trigger on \emph{local} + \emph{specific}. The Gate can be crossed proactively by the writer or automatically when the system detects the threshold is exceeded.


\subsubsection{Negotiation}
\label{sec:negotiation}

After crossing the Gate, negotiation targets the team's shared understanding. The following dimensions define the Negotiation design space:

\paragraph{Propose --- initiating a BNA modification.}

\begin{table}[h]
\small
\begin{tabular}{@{}lp{2cm}p{4.5cm}@{}}
\toprule
\textbf{Dimension} & \textbf{Design Question} & \textbf{Option Space} \\
\midrule
Initiation & Who/what initiates & Writer, system-suggested, owner-only \\
Scope & What can be proposed & Modify intent, add/delete intent, adjust dependency, reassign \\
Framing & What context accompanies & Impact analysis, affected parties, writing preview, decision history \\
\bottomrule
\end{tabular}
\end{table}

\paragraph{Deliberate --- affected parties review.}

\begin{table}[h]
\small
\begin{tabular}{@{}lp{2cm}p{4.5cm}@{}}
\toprule
\textbf{Dimension} & \textbf{Design Question} & \textbf{Option Space} \\
\midrule
Participation & Who participates & Direct dependents, all owners, all members \\
Mechanism & How they deliberate & Voting, threaded discussion, synchronous meeting, silent approval (timeout = consent) \\
Visibility & What participants see & Impact on their intents/writing, others' opinions, similar past decisions \\
Response & How to express opinion & Approve, reject, counter-propose, defer, request more info \\
\bottomrule
\end{tabular}
\end{table}

\paragraph{Resolve --- reaching a decision.}

\begin{table}[h]
\small
\begin{tabular}{@{}lp{2cm}p{4.5cm}@{}}
\toprule
\textbf{Dimension} & \textbf{Design Question} & \textbf{Option Space} \\
\midrule
Criteria & What counts as decided & Unanimous, majority, stakeholder agreement, proposer decides \\
Fallback & If not resolved & Maintain status quo, escalate, auto-resolve after timeout, split into sub-decisions \\
Outcome & How it takes effect & Update BNA, notify affected parties, record decision history, trigger awareness recomputation \\
\bottomrule
\end{tabular}
\end{table}


\subsubsection{Defining a Coordination Pipeline with the Design Space}
\label{sec:pipeline-examples}

These dimensions define the coordination design space for collaborative writing. Teams select values across these dimensions to define their own coordination pipeline. Two examples show how the same design space supports fundamentally different coordination approaches.

\paragraph{Example 1: Two co-authors writing a blog post synchronously.} Alice and Bob brainstorm face-to-face with a simple living outline as BNA. They configure: Trigger = ambient; Gate = fundamental + team-wide only; Group Interest = default won't-care; Mechanism = synchronous discussion; Criteria = proposer-decides. Lightweight awareness, loose gate, instant negotiation. Most changes resolve in personal space; only directional shifts cross to team space.

\paragraph{Example 2: Five co-authors writing an academic paper asynchronously.} Distributed across three time zones, with a structured living outline and explicit dependencies. They configure: Trigger = post-action; full relationship tracking including Writing $\leftrightarrow$ Others; Gate = local + specific triggers; Group Interest = default need-to-know; Mechanism = voting with 48-hour deadline; Participation = direct dependents; Criteria = stakeholder agreement; Fallback = timeout maintains status quo, escalate to all. Fine-grained awareness, strict gate, formal voting. The system proactively detects drift and suggests initiating negotiation.

\paragraph{Awareness is continuous.} Awareness operates not only before the Gate. The dimensions defined in \S\ref{sec:awareness}---relationship changes, impact scope, group interest---are needed at every stage of the coordination loop: during writing, when exploring potential changes, during Deliberate (participants need to understand how a proposal affects their writing), and after Resolve (re-sensing the updated relationship state). Before the Gate, awareness serves the question ``should I cross to team space?''; after the Gate, it serves cognition and decision-making within the negotiation process.

The design space is AI-agnostic: without AI, teams can manually check writing-intent relationships, manually judge impact, and negotiate face-to-face; with AI, the cost of perception and judgment at each dimension drops substantially, but the framework itself does not depend on AI.

\section{\system{}: An Open Platform for Coordination}
\label{sec:system}

Section~\ref{sec:framework} defined a coordination design space for collaborative writing---Awareness dimensions, a Gate threshold, and Negotiation dimensions. This section introduces \system{}, an open platform built on that design space. Rather than a fixed implementation, \system{} provides an extensible architecture: developers contribute new awareness and negotiation capabilities through declarative protocols; teams select from these capabilities and configure their own coordination pipeline. We describe the platform in three layers: Foundation (\S\ref{sec:foundation}), Protocol (\S\ref{sec:protocol}), and Team Interface (\S\ref{sec:team-interface}).

% ─── FIGURE: System Architecture ───
\begin{figure*}[t]
\centering
\fbox{\parbox{0.95\textwidth}{\centering\vspace{3cm}
\textbf{Figure: \system{} Architecture}\\[0.5em]
\small Three layers: Foundation (BNA as Living Outline + Data Model + View Layer), Protocol (Awareness Protocol + Negotiation Protocol), Team Interface (Setup + Writing + Coordination). Arrows show data flow: Data Model feeds into Protocol functions; Protocol outputs render through View Layer; Team Interface configures which functions and paths are active.
\vspace{3cm}}}
\caption{\system{} architecture. The Foundation layer provides the living outline (BNA), abstract data model, and view layer with entity-based UI slots. The Protocol layer enables developers to contribute awareness functions and negotiation paths through declarative registration. The Team Interface lets teams configure their coordination pipeline through guided setup. All built-in capabilities use the same protocol as third-party extensions.}
\label{fig:architecture}
\end{figure*}


\subsection{Foundation}
\label{sec:foundation}

\subsubsection{Living Outline as BNA}

\system{} adopts the \emph{living outline} as its BNA: a structured, continuously evolving outline where each node carries an intent description, dependency annotations, and ownership assignments. The living outline satisfies conditions C1--C5 from Section~\ref{sec:bna}:

\begin{itemize}[leftmargin=1.5em, itemsep=1pt]
\item \textbf{C1 (Intent units):} Each outline node describes an intent---more abstract than prose, more concrete than a topic heading.
\item \textbf{C2 (Dependencies):} The outline supports typed dependencies between nodes (depends-on, supports, must-be-consistent) and hierarchical structure (sections contain sub-intents).
\item \textbf{C3 (Ownership):} Each node has a designated owner.
\item \textbf{C4 (Traceability):} The system continuously computes the correspondence between writing and intent.
\item \textbf{C5 (Negotiability):} The outline is the site of structured proposals---any modification to intents, dependencies, or assignments flows through the coordination pipeline.
\end{itemize}

\noindent The process of maintaining the living outline during writing itself pushes the negotiation boundary: writers continuously update their understanding of intent, and the outline continuously evolves to reflect the team's latest consensus.


\subsubsection{Data Model}
\label{sec:data-model}

\system{} abstracts the negotiation space into a unified data model with two layers: a \textbf{structure layer} (the document's current state) and an \textbf{interaction layer} (the coordination process's dynamic state). This data model is the sole input interface for all functions---a function sees the complete current state and interaction history.

\paragraph{Structure layer --- BNA + Writing.} The structure layer directly corresponds to the negotiation space defined in Section~\ref{sec:bna}:

\begin{itemize}[leftmargin=1.5em, itemsep=1pt]
\item \textbf{IntentItem}: An intent unit---id, content description, hierarchical position (\texttt{parentId}), creator, last modifier, change status (stable / proposed / disputed).
\item \textbf{Dependency}: A relationship between two IntentItems---type (depends-on / must-be-consistent / builds-upon), directionality, source (manual / AI-suggested / AI-confirmed).
\item \textbf{Section}: A writing unit bound to an IntentItem via \texttt{intentId}---assignee, content, and per-paragraph attribution (who wrote each paragraph, who last edited it, when).
\end{itemize}

\noindent IntentItem and Section are bound through \texttt{section.intentId}---the concrete expression of the BNA $\rightarrow$ Writing relationship. The BNA layer is replaceable: in a living outline, IntentItem is a hierarchical node; in an argument map, it is a claim/evidence node; but the interface remains the same.

\paragraph{Interaction layer --- dynamic coordination state.} Beyond document structure, functions need access to the coordination process itself---what previous functions have output, whether proposals are in progress, and what past resolutions have established:

\begin{itemize}[leftmargin=1.5em, itemsep=1pt]
\item \textbf{FunctionResult}: Output from a previously run function---which function, which target, output data, timestamp. Subsequent functions can read prior results (e.g., \texttt{assess-change-impact} reads \texttt{check-drift}'s output without recomputing).
\item \textbf{Proposal}: An active or resolved proposal---target IntentItem, proposed change type, proposer, status, votes, attached function results, resolution timestamp.
\end{itemize}

\noindent The system assembles the complete data model before invoking any function---developers do not need to understand how data is constructed from the real-time collaboration state.


\subsubsection{View Layer}
\label{sec:view-layer}

The View Layer exposes declarative rendering interfaces on top of the data model. Developers declare ``show this type of UI on this data object''---the system handles rendering.

The negotiation space naturally produces two core views and a global panel:

\begin{itemize}[leftmargin=1.5em, itemsep=1pt]
\item \textbf{Outline View} (left): Renders the BNA. Currently a living outline; replacing the BNA form changes this view (e.g., to an argument map visualization). Developer declarations remain the same---the BNA adapter handles rendering.
\item \textbf{Writing View} (right): Renders writing content. A paragraph editor bound to IntentItems in the Outline View.
\item \textbf{Panel}: Global display space with pre-defined slots for analysis results, diffs, voting interfaces, and cross-object summaries.
\end{itemize}

Each entity type offers UI capabilities determined by its nature. The system provides base-level rendering capabilities; complex components (voting widgets, discussion threads) are built on top.

\begin{table}[t]
\small
\caption{View Layer: entity-based UI capabilities and their slots.}
\label{tab:view-layer}
\begin{tabular}{@{}llp{2cm}p{3.5cm}@{}}
\toprule
\textbf{Entity} & \textbf{Capability} & \textbf{Slots} & \textbf{Description} \\
\midrule
\multirow{3}{*}{IntentItem} & Indicator & icon, value, color & Status marks beside the item \\
& Expandable & content & Expandable area below the item \\
& Actions & buttons, menu & Action entry points on the item \\
\midrule
\multirow{3}{*}{Paragraph} & Inline & range, style, tooltip & In-text: highlights, underlines, marks \\
& Side & anchor, content & Beside text: comments, annotations \\
& Block & position, content & Above/below: banners, alerts \\
\midrule
Dependency & State & style, color, label & Connection visual state \\
\midrule
\multirow{2}{*}{Panel} & Card & title, content, footer & Generic card container \\
& Widget & content & Free rendering area \\
\bottomrule
\end{tabular}
\end{table}

\noindent Based on these base capabilities, we build higher-level components (ResultList, DiffView, VoteWidget, CommentThread) as rendering implementations for built-in functions. Developers can build custom components using the same base capabilities.


\subsection{Protocol}
\label{sec:protocol}

\system{} enables extensibility through two declarative protocols: an \textbf{Awareness Protocol} for contributing sensing capabilities (corresponding to \S\ref{sec:awareness}), and a \textbf{Negotiation Protocol} for contributing coordination paths (corresponding to \S\ref{sec:negotiation}).

\subsubsection{Awareness Protocol}

An awareness function declares four fields:

\begin{description}[leftmargin=1.5em, nosep, itemsep=3pt]
\item[Target:] What position in the negotiation space---IntentItem, Section, Dependency, or Document.
\item[Trigger:] When to run---ambient (continuous), on-action (event-triggered), on-demand (manual), post-resolve (after resolution).
\item[Executor:] How to detect---\texttt{prompt} (AI model + configuration) or \texttt{local} (local function + configuration). Detection logic and configurable options (e.g., focus areas) are defined within the executor; \texttt{fill}-type config fields let teams input custom prompt text that is injected into the AI call.
\item[Display:] Where to render results---entity type + base capability + slot data bindings using template syntax (e.g., \texttt{\{\{output.coverage\}\}}).
\end{description}


\subsubsection{Negotiation Protocol}

A negotiation path declares fields corresponding to Propose $\rightarrow$ Deliberate $\rightarrow$ Resolve:

\begin{table}[h]
\small
\begin{tabular}{@{}llp{5cm}@{}}
\toprule
\textbf{Stage} & \textbf{Field} & \textbf{Options} \\
\midrule
\multirow{3}{*}{Propose} & Initiation & Writer, system-suggested, owner-only \\
& Scope & Modify/add/delete intent, adjust deps, reassign \\
& Framing & Auto-attach which awareness results as context \\
\midrule
\multirow{3}{*}{Deliberate} & Participation & Direct dependents, all owners, all members \\
& Mechanism & Voting, discussion, silent-approval \\
& Response & Approve, reject, counter-propose, defer \\
\midrule
\multirow{3}{*}{Resolve} & Criteria & Unanimous, majority, stakeholder, proposer \\
& Fallback & Status quo, timeout auto-resolve, escalate \\
& Outcome & Update BNA, notify, record history, re-trigger awareness \\
\bottomrule
\end{tabular}
\end{table}


\subsubsection{How It Works}

Registering a function via \texttt{registerFunction()} integrates it into the system like installing a plugin:

\begin{enumerate}[leftmargin=1.5em, itemsep=1pt]
\item \textbf{Register}: The function declares target, trigger, executor, output schema, and display bindings.
\item \textbf{Trigger}: The system invokes the function at the right moment---ambient functions run on document changes; on-action functions run on specific events (e.g., proposal submission).
\item \textbf{Execute}: The system passes the current data model (IntentItems + Sections + Dependencies + interaction layer). For prompt executors, data and team configuration are injected into the prompt; for local executors, the function runs directly.
\item \textbf{Render}: The system binds output data to the declared entity slots---no UI code needed from the developer.
\item \textbf{Store}: Output is stored as a FunctionResult in the interaction layer, available to subsequent functions.
\end{enumerate}

\noindent Negotiation path registration works similarly: declaring Propose/Deliberate/Resolve configuration automatically provisions the corresponding UI (proposal forms, voting interfaces, discussion threads) and resolution logic.


\subsubsection{Examples}

Figure~\ref{fig:protocol-examples} shows two complete function declarations.

% ─── FIGURE: Protocol Examples ───
\begin{figure*}[t]
\centering
\fbox{\parbox{0.95\textwidth}{\centering\vspace{4cm}
\textbf{Figure: Two Function Declarations}\\[0.5em]
\small
\textbf{Left: check-drift (AI-powered).} Target: Section. Trigger: ambient. Executor: prompt with \texttt{focusPrompt} fill-config. Output: coverage score + per-sentence alignment array. Display: Paragraph.Inline (sentence highlights by alignment) + IntentItem.Indicator (coverage value).\\[1em]
\textbf{Right: check-heading-consistency (local).} Target: Section. Trigger: ambient. Executor: local function comparing heading text to intent content. Output: heading, intentContent, consistent boolean. Display: IntentItem.Indicator (warning icon when inconsistent) + Paragraph.Side (mismatch note).
\vspace{4cm}}}
\caption{Two function declarations through the Awareness Protocol. Left: an AI-powered drift detector that highlights sentence-level alignment in Writing View and shows coverage in Outline View. Right: a local consistency checker requiring no AI. Both use the same protocol, the same data model, and the same view layer slots.}
\label{fig:protocol-examples}
\end{figure*}


\subsubsection{Built-in Capabilities}
\label{sec:built-in}

Based on the design space in Section~\ref{sec:framework}, we pre-define capabilities for each stage of the coordination loop. These are registered through the same protocol as third-party extensions---we are the first users of our own protocol, not a privileged layer.

\paragraph{Awareness Functions} are cognitive support capabilities that help writers and teams sense, understand, and evaluate relationship changes in the negotiation space. They are not limited to the pre-Gate phase---at any point in the coordination loop where understanding relationship state is needed, an awareness function can be triggered. A function can be mounted on writing actions (editing an intent, completing a paragraph) or on Negotiation Path nodes (assisting participants during Deliberate).

We organize our built-in functions by the cognitive need they serve:

\textbf{Sensing relationship state}---helping writers continuously understand their position in the negotiation space:
\begin{itemize}[leftmargin=1.5em, itemsep=1pt]
\item \textbf{check-drift} (Writing $\leftrightarrow$ Intent, Writing $\leftrightarrow$ BNA): Continuously detects whether writing covers intent and aligns with team direction. Highlights drifted sentences in Writing View, shows coverage in Outline View. The most fundamental cognitive need: \emph{what is the relationship between what I wrote and what the team agreed?}
\item \textbf{check-cross-consistency} (Writing $\leftrightarrow$ Others): Detects whether my writing is consistent with others' writing through dependency chains. When Alice revises the research question, this function alerts Bob that his Method section may need adjustment.
\item \textbf{discover-dependencies} (BNA internal): Analyzes outline structure to discover and suggest implicit dependencies between intent units. Runs during Setup and after outline changes.
\end{itemize}

\textbf{Assessing impact}---helping writers understand the consequences of a potential change, providing data for Gate threshold decisions:
\begin{itemize}[leftmargin=1.5em, itemsep=1pt]
\item \textbf{assess-change-impact} (Scope, Severity, Propagation): Evaluates the impact of modifying or deleting an intent---which sections are affected, how significant the impact is, how far it propagates through dependency chains. Output directly feeds Gate threshold evaluation.
\item \textbf{preview-change-on-writing} (Propagation): Projects a potential modification onto affected sections, showing a diff of how existing writing would change. Writers use it to evaluate ``should I propose this?''; participants in Deliberate use it to understand ``what does this proposal mean for my writing?''---the same function serving different cognitive needs at different stages.
\end{itemize}

\textbf{Supporting the negotiation process}---providing cognitive support during Propose, Deliberate, and Resolve:
\begin{itemize}[leftmargin=1.5em, itemsep=1pt]
\item \textbf{suggest-intent-for-orphan}: When a writer produces content outside any intent, suggests creating a new IntentItem.
\item \textbf{frame-proposal}: Automatically organizes proposal context---attaches impact analysis, links to related past resolutions, lists affected parties.
\item \textbf{preview-alternatives}: During Deliberate, compares multiple proposed alternatives side-by-side showing their impact on writing.
\item \textbf{preview-resolution-effect}: During Resolve, previews the final state of BNA and writing if the proposal is accepted.
\end{itemize}


\paragraph{Negotiation Paths} define the procedural rules of negotiation---who participates, how they deliberate, what counts as resolution. Each path comes with default-mounted Awareness Functions that provide cognitive support at key nodes. Teams can enable or disable these defaults to prevent excessive detection.

\begin{table}[t]
\small
\caption{Built-in Negotiation Paths with default-mounted Awareness Functions.}
\label{tab:negotiation-paths}
\begin{tabular}{@{}lllp{3.5cm}@{}}
\toprule
\textbf{Path} & \textbf{Mechanism} & \textbf{Criteria} & \textbf{Default Functions} \\
\midrule
Inform & Notify & Immediate & assess-change-impact \\
Ask for Input & Solicit & Single-approval & assess-change-impact, preview-change-on-writing \\
Team Vote & Voting & Majority & frame-proposal, preview-change-on-writing, preview-resolution-effect \\
Discussion & Threaded & Proposer-closes & frame-proposal, preview-alternatives \\
\bottomrule
\end{tabular}
\end{table}

\noindent For example, the Team Vote path defaults to running \texttt{frame-proposal} at Propose (auto-organizing context), \texttt{preview-change-on-writing} at Deliberate (showing voters the impact), and \texttt{preview-resolution-effect} before Resolve (letting the team preview the outcome). Teams can disable any of these---a small team might turn off \texttt{preview-resolution-effect} as unnecessary overhead.

All capabilities are registered through the protocol, not hardcoded. As the community contributes new functions (terminology consistency checkers, citation validators, formatting harmonizers) and new paths, teams select from an expanding pool.


\subsection{Team Interface}
\label{sec:team-interface}

\subsubsection{Setup: Defining the Coordination Pipeline}
\label{sec:setup}

Before writing begins, teams define their coordination pipeline through a guided Setup interface. This process resembles teams authoring their own MetaRule---selecting which capabilities to activate, configuring Gate thresholds, and choosing Negotiation Paths.

The Setup interface presents questions organized around the pre-defined capabilities:

\begin{itemize}[leftmargin=1.5em, itemsep=1pt]
\item \textbf{Awareness}: Enable check-drift? If so, what should the checker focus on? (configures \texttt{focusPrompt}). Enable cross-consistency checking? Enable dependency discovery?
\item \textbf{Gate}: What level of drift should notify the team? (configures Severity threshold). Auto-trigger when others are affected, or let writers decide? (configures Scope threshold + Agency).
\item \textbf{Negotiation}: How should intent modifications be decided? (selects Path). Who participates? (configures Participation). How long before timeout? (configures deadline).
\end{itemize}

\noindent After answering, the system generates a visual coordination pipeline---the complete Awareness $\rightarrow$ Gate $\rightarrow$ Negotiation flow, showing which functions and paths are active at each node. Most teams can start writing immediately. Teams wanting finer control can click any node in the pipeline to adjust configuration.

% ─── FIGURE: Setup ───
\begin{figure}[t]
\centering
\fbox{\parbox{0.95\columnwidth}{\centering\vspace{2.5cm}
\textbf{Figure: Setup Interface}\\[0.5em]
\small Left: guided questions organized by Awareness / Gate / Negotiation. Right: generated coordination pipeline (editable). Nodes show active functions and paths.
\vspace{2.5cm}}}
\caption{The Setup interface. Teams answer questions about their collaboration context; the system generates a visual coordination pipeline that can be further customized.}
\label{fig:setup}
\end{figure}


\subsubsection{Writing and Coordination}

Once configured, all interfaces load adaptively based on the team's pipeline: which functions are active determines what annotations appear in Outline View and Writing View; which path is selected determines what interaction the coordination interface provides (voting, discussion, silent approval). The interface the team sees is the direct embodiment of the pipeline they defined in Setup.


\subsubsection{Configuration Examples}
\label{sec:config-examples}

Two examples show how different teams use Setup to define fundamentally different coordination pipelines (Figure~\ref{fig:pipeline-comparison}).

\paragraph{Example 1: Two co-authors, synchronous blog post.} Setup: enable check-drift (focusPrompt: ``check overall argument flow''), disable cross-consistency; Gate: fundamental severity only; Path: Inform; Participation: all.

\noindent \textbf{Generated pipeline}: Writing $\rightarrow$ check-drift (ambient, loose) $\rightarrow$ Gate: severity $\geq$ fundamental $\rightarrow$ Inform (notify, no response needed). \emph{Effect}: Lightweight awareness, loose gate. Outline View shows basic coverage; Writing View highlights only directional drift. Most changes resolve in personal space through conversation.

\paragraph{Example 2: Five co-authors, asynchronous academic paper.} Setup: enable check-drift (focusPrompt: ``check argument structure and evidence alignment''), enable cross-consistency and dependency discovery; Gate: local + specific triggers; Path: Team Vote (majority, 48h); Participation: direct dependents.

\noindent \textbf{Generated pipeline}: Writing $\rightarrow$ check-drift (strict) + check-cross-consistency $\rightarrow$ Gate: severity $\geq$ local AND scope $\geq$ specific $\rightarrow$ Team Vote [Propose: frame-proposal; Deliberate: preview-change-on-writing + preview-alternatives; Resolve: majority, 48h, preview-resolution-effect] $\rightarrow$ Update BNA $\rightarrow$ re-trigger awareness. \emph{Effect}: Fine-grained awareness, strict gate, formal voting. Full annotation in both views; Deliberate shows diffs and alternative comparisons automatically.

% ─── FIGURE: Pipeline Comparison ───
\begin{figure*}[t]
\centering
\fbox{\parbox{0.95\textwidth}{\centering\vspace{3cm}
\textbf{Figure: Two Team Pipelines}\\[0.5em]
\small Left: lightweight pipeline (two-person blog)---two nodes (check-drift $\rightarrow$ Inform), minimal annotations. Right: full pipeline (five-person paper)---multiple nodes with branching (check-drift + cross-consistency $\rightarrow$ strict Gate $\rightarrow$ Team Vote with frame-proposal, preview, and resolution preview), rich annotations in both views.
\vspace{3cm}}}
\caption{The same design space produces fundamentally different coordination pipelines. Left: a two-person synchronous team with lightweight awareness and informal negotiation. Right: a five-person distributed team with fine-grained awareness, strict gate thresholds, and formal voting with AI-assisted cognitive support at each negotiation stage.}
\label{fig:pipeline-comparison}
\end{figure*}


\bibliographystyle{ACM-Reference-Format}
\bibliography{references}

\end{document}
